\scanpage{76}\mspage{37}
% slanted
\margin{\deleted{\unsure{inutile} ?}}
il fixe aussi $-\deleted{\unsure{x|\xi}}\unsure{s|\xi|x}$.] S'il fixe $\deleted{\unsure{s|\xi}}x$ et $y$ tels que $x\wedge y\in i=\{e,-e\}\subset\deleted{\unsure{\det}}\overrightarrow{\unsure{\pi|\Pi}}\deleted{\unsure{{}^*}}=(\det M)^*$, alors il fixe aussi $\inserted{z=}-(x+y)$, \inserted{donc aussi $x+y$.} Cela nous permet de nous propager à partir d'un \inserted{$i$-}repère $(x,y)$, \deleted{fixé par $u$} \deleted{\unsure{on le transforme soit par $(x,y+x$|les remplaçant par $(x,y+x)$}} i.e. à partir d'une $i$-face $(x,y,z)$, en faisant soit \deleted{des transpositions}\inserted{une permutation} sur $(x,y,z)$ --- p. ex. les deux transpositions entre $x$ et $y$, et entre $y$ et $z$ --- et la transformation $(x,y,z)\longmapsto(x,-y,y-x)$ ---
\margin{\unsure{1}}
en termes matriciels, avec tout \inserted{$i$-}repère \uline{de $M$} $r=(x,y)$, interprété comme isom
\[
 r:\Lambda^2\xrightarrow{\sim}M
\]
qui soit formé de ``composantes'' $x,y$ fixées par $u$ (on dit que $r$ est ``$u$-fixe''), il en est de \unsure{$=$|même} de $r\circ\tau_0^{\unsure{r|v}},\ r\circ\tau_1^r,\ r\circ\tau_\infty^r$ (cf (82) \inserted{p. 15'}) .
Cela signifie que si $r$ est $u$-fixe, de $=$ $r\circ g$, où
\[
 \bigl(\qquad g\in\underset{\substack{\Vert\\\operatorname{Aut}_\Lambda(\Lambda^2)}}{\mathrm{\deleted{\unsure{S}}Gl}(2,\Lambda)}
 \quad\begin{gathered}\text{est dans le sous-groupe}\\\text{engendré par }\tau_0^r,\tau_1^r,\tau_\infty^r.\end{gathered}
\]
Or ce sous-groupe est l'image, par
\[
 \mathrm{Sl}(2,\mathbb Z)\longrightarrow\mathrm{Sl}(2,\Lambda)\subset\mathrm{Gl}(2,\Lambda)
\]
du sous-groupe de $\mathrm{Sl}(2,\mathbb Z)$ engendré par ces $=$s matrices $\tau_i^r$ (interprétées comme étant à
% archive sheet number: 75
